\section[\appendixname~\thesection]{Roughness at the daily frequency}\label{app:rough}

The increment-scaling estimate of the Hurst exponent of daily log Parkinson
variance averages 0.060 across stocks (10th--90th percentile 0.053--0.069).
Read literally, this indicates volatility much rougher than Brownian motion,
for which the exponent is 0.5. We ask whether a
process without fractional roughness, observed through the Parkinson
estimator, produces the same value. Latent log variance is the sum of two
independent AR(1) components, a Markovian model. The observed series adds the
sampling error of the range estimator, $\eta_t = \ln(R_t^2/4\ln 2)$ minus its
mean, where $R_t$ is the range of a standard Brownian motion over one day
sampled at 390 one-minute steps; its variance, 0.36, is fixed by theory and
is about a third of the variance of observed log variance. The two
autoregressive coefficients and the variance split are fitted to the median
autocorrelations of observed log variance at nine lags from 1 to 250 days,
which gives a fast component with coefficient 0.66 and a slow one with
coefficient 0.995. Over 200 simulated paths of 6136 days the same estimator
gives 0.064 (standard deviation 0.004) on the observed series, matching the
data, and 0.19 on the noise-free latent series. Two features of daily data
therefore produce the low estimate without any roughness: the sampling noise
of the range, which lowers the estimate from 0.19 to 0.06, and a fast
mean-reverting component, which already holds it well below 0.5 at lags of up
to 21 days. Daily range data cannot identify roughness, in line with
\citet{cont2024rough}. The code is module 16 of the repository.

\section[\appendixname~\thesection]{Additional forecast results}\label{app:extra}

\textit{QLIKE.} Variance forecasts are obtained from the log forecasts with a
point-in-time smearing correction: the mean of the exponentiated past
out-of-sample errors whose targets had been observed by the origin, with at
least 26 such errors. Against HAR-X, no specification has a significantly
lower QLIKE loss at the 5\% level; the largest Diebold--Mariano statistic in
favor of persistence is 1.67 (cross-sectional extension, one day), and model
$C$ is significantly worse at one month ($-2.08$).

\textit{Machine learning.} The lasso, ridge and elastic-net penalties are
chosen by five-fold expanding time-series cross-validation with a gap of
$\lceil h/5\rceil$ rows between training and validation folds. Gradient
boosting (LightGBM, learning rate 0.05) chooses the number of leaves
$\{2,4,15\}$, the minimum leaf size $\{20,50\}$ and the number of trees
$\{25,\dots,800\}$ by three-fold time-series cross-validation in each
training window. The random forest uses 200 trees with a minimum leaf of 20.
Each learner is refit every 20 weeks on the embargoed training rows. With
fixed defaults (400 trees, 31 leaves) gradient boosting loses 15.9\%,
11.9\% and 21.2\% to HAR at the three horizons; tuned, it loses 2.9\%, 0.2\%
and 8.3\%.

\section[\appendixname~\thesection]{The market-level gain}\label{app:market}

On the market-level series the persistence state improves on market HAR-X by
3.4\% and 2.0\% at one day and one week (Table~\ref{tab:designs}). An
exploratory analysis, specified after that result and therefore not a formal
test, asks what accounts for the gain. In a specification in levels, the
interaction of the VIX with the logarithm of its maximum over the previous
three years has a negative coefficient ($t=-2.7$ at both horizons): the
slope of future realized variance on the VIX is about 40\% lower when a
crisis is recent. In logs the interaction disappears ($t=0.33$ and $0.26$),
and the three-year maximum enters only as a level shift, which is absorbed
by the trailing 250-day gap between VIX-implied and realized variance. That
gap alone improves on the log market HAR-X model by 5.0\% to 6.3\%
(Diebold--Mariano statistics of 2.7 to 3.5). The market-level gain is
therefore best interpreted as the slowly decaying variance risk premium after
crises \citep{bates2000post, andersen2015risk}, for which the persistence
state may act as a proxy because both depend on how recently the last crisis
occurred. A test of this interpretation requires data not used here, such as
earlier index history or other markets.

\section[\appendixname~\thesection]{Registration record and audit}\label{app:ledger}

The study was carried out in stages. Before each stage the specifications,
comparisons, tests and decision rules were written into an experiment ledger
kept under version control; results were added afterwards and earlier entries
were not edited. The ledger, the code and a list of every correction are in
the repository named in the Data Availability Statement. The main stages
were:
\begin{itemize}
\item an embargo on training observations whose targets were not yet
realized at the origin, which removed an apparent gain of the persistence
model at the monthly horizon;
\item three HAR-X extensions with a fixed decision rule and Holm adjustment;
\item pooled, term-structure and market-level designs, point-in-time
winsorization and real-time portfolio normalization;
\item a full audit of the pipeline, whose corrections were registered before
the final run: aligned information timing (Table~\ref{tab:timing}), missing
rather than floored zero-range days, memory estimated on log variance
throughout, a Newey--West bandwidth that grows with the sample, point-in-time
QLIKE, time-series cross-validation for every learner, a level-corrected
GARCH benchmark, common evaluation cells and dates, and portfolio returns in
excess of the bill rate;
\item the window-inclusion analysis (Table~\ref{tab:window}), the
Giacomini--White tests and the timing comparison, each registered with a
prediction before it was run.
\end{itemize}
% CO-AUTHORS: decide whether to add one sentence noting that an earlier
% version (arXiv:2605.24285) reported larger gains, and that the
% differences come from the corrections listed above.
